\documentclass[10pt,a4paper]{amsart}
\usepackage{amsmath,amssymb,amsthm,graphicx}
% Retained minimal preamble for the Stacks Project chapter more-algebra.tex (More on Algebra).
% The chapter file pinned-original-source.tex carries only \input{preamble} on its first line, so no
% preamble text lives inside the chapter itself.  The Stacks Project uses one single shared preamble for
% every chapter; the snapshot stages it identically as sources/stacks-chow/preamble.tex and
% sources/stacks-etale-cohomology/preamble.tex.  The block below is copied mechanically, line for line,
% from sources/stacks-chow/preamble.tex lines 14-18 (the Xy-pic setup) and 154-217 (the native theorem
% declarations together with the mathematical macros).  Nothing was added, removed or redefined; the
% wrapper class and the amsmath/amssymb/amsthm/graphicx base packages are supplied by the assembler.
% Chapter: more-algebra.tex -- More on Algebra.
% For commutative diagrams we use Xy-pic
\usepackage[all]{xy}

% We use 2cell for 2-commutative diagrams.
\xyoption{2cell}

% Theorem environments.
%
\theoremstyle{plain}
\newtheorem{theorem}[subsection]{Theorem}
\newtheorem{proposition}[subsection]{Proposition}
\newtheorem{lemma}[subsection]{Lemma}

\theoremstyle{definition}
\newtheorem{definition}[subsection]{Definition}
\newtheorem{example}[subsection]{Example}
\newtheorem{exercise}[subsection]{Exercise}
\newtheorem{situation}[subsection]{Situation}

\theoremstyle{remark}
\newtheorem{remark}[subsection]{Remark}
\newtheorem{remarks}[subsection]{Remarks}

\numberwithin{equation}{subsection}

% Macros
%
\def\lim{\mathop{\mathrm{lim}}\nolimits}
\def\colim{\mathop{\mathrm{colim}}\nolimits}
\def\Spec{\mathop{\mathrm{Spec}}}
\def\Hom{\mathop{\mathrm{Hom}}\nolimits}
\def\Ext{\mathop{\mathrm{Ext}}\nolimits}
\def\SheafHom{\mathop{\mathcal{H}\!\mathit{om}}\nolimits}
\def\SheafExt{\mathop{\mathcal{E}\!\mathit{xt}}\nolimits}
\def\Sch{\mathit{Sch}}
\def\Mor{\mathop{\mathrm{Mor}}\nolimits}
\def\Ob{\mathop{\mathrm{Ob}}\nolimits}
\def\Sh{\mathop{\mathit{Sh}}\nolimits}
\def\NL{\mathop{N\!L}\nolimits}
\def\CH{\mathop{\mathrm{CH}}\nolimits}
\def\proetale{{pro\text{-}\acute{e}tale}}
\def\etale{{\acute{e}tale}}
\def\QCoh{\mathit{QCoh}}
\def\Ker{\mathop{\mathrm{Ker}}}
\def\Im{\mathop{\mathrm{Im}}}
\def\Coker{\mathop{\mathrm{Coker}}}
\def\Coim{\mathop{\mathrm{Coim}}}

% Boxtimes
%
\DeclareMathSymbol{\boxtimes}{\mathbin}{AMSa}{"02}

%
% Macros for moduli stacks/spaces
%
\def\QCohstack{\mathcal{QC}\!\mathit{oh}}
\def\Cohstack{\mathcal{C}\!\mathit{oh}}
\def\Spacesstack{\mathcal{S}\!\mathit{paces}}
\def\Quotfunctor{\mathrm{Quot}}
\def\Hilbfunctor{\mathrm{Hilb}}
\def\Curvesstack{\mathcal{C}\!\mathit{urves}}
\def\Polarizedstack{\mathcal{P}\!\mathit{olarized}}
\def\Complexesstack{\mathcal{C}\!\mathit{omplexes}}
% \Pic is the operator that assigns to X its picard group, usage \Pic(X)
% \Picardstack_{X/B} denotes the Picard stack of X over B
% \Picardfunctor_{X/B} denotes the Picard functor of X over B
\def\Pic{\mathop{\mathrm{Pic}}\nolimits}
\def\Picardstack{\mathcal{P}\!\mathit{ic}}
\def\Picardfunctor{\mathrm{Pic}}
\def\Deformationcategory{\mathcal{D}\!\mathit{ef}}

% --- appended 2026-10-12 (Chief): the project own definitions of the reference/slogan/history
% display environments, copied mechanically from the shared Stacks preamble lines 8-12
% (\usepackage{verbatim} plus the three \newenvironment lines, which the chapter preamble keeps
% verbatim from the project preamble). Nothing added, removed or redefined: these environments
% carry no mathematical content and the project itself typesets them as comments. Without these
% three lines a counted statement whose own text uses them fails with Environment slogan undefined.
\usepackage{verbatim}
\newenvironment{reference}{\comment}{\endcomment}
\newenvironment{slogan}{\comment}{\endcomment}
\newenvironment{history}{\comment}{\endcomment}
\usepackage{cleveref}
\crefname{originallemma}{original lemma}{original lemmas}
\Crefname{originallemma}{Original lemma}{Original lemmas}
\makeatletter
\newlabel{algebra-lemma-characterize-zero-local}{{00HN}{source}{}{sourcepointer-1}{}}
\newlabel{algebra-lemma-characterize-zero-local@cref}{{[originallemma][0][]00HN}{[1][0][]source}{}{}{}}
\makeatother
\crefname{originallemma}{original lemma}{original lemmas}
\Crefname{originallemma}{Original lemma}{Original lemmas}
\makeatletter
\newlabel{algebra-lemma-hom-from-finitely-presented}{{0583}{source}{}{sourcepointer-2}{}}
\newlabel{algebra-lemma-hom-from-finitely-presented@cref}{{[originallemma][0][]0583}{[1][0][]source}{}{}{}}
\makeatother
\crefname{originallemma}{original lemma}{original lemmas}
\Crefname{originallemma}{Original lemma}{Original lemmas}
\makeatletter
\newlabel{lemma-characterize-PD-modules}{{37586}{source}{}{sourcepointer-3}{}}
\newlabel{lemma-characterize-PD-modules@cref}{{[originallemma][0][]37586}{[1][0][]source}{}{}{}}
\makeatother
\crefname{originallemma}{original lemma}{original lemmas}
\Crefname{originallemma}{Original lemma}{Original lemmas}
\makeatletter
\newlabel{lemma-generalized-valuation-ring}{{37638}{source}{}{sourcepointer-4}{}}
\newlabel{lemma-generalized-valuation-ring@cref}{{[originallemma][0][]37638}{[1][0][]source}{}{}{}}
\makeatother
\crefname{originallemma}{original lemma}{original lemmas}
\Crefname{originallemma}{Original lemma}{Original lemmas}
\makeatletter
\newlabel{lemma-generalized-valuation-ring-modules}{{37675}{source}{}{sourcepointer-5}{}}
\newlabel{lemma-generalized-valuation-ring-modules@cref}{{[originallemma][0][]37675}{[1][0][]source}{}{}{}}
\makeatother
\title{Over a ring whose localizations satisfy the Warfield conditions every finitely presented module is a summand of a finite sum of quotients by principal ideals}
\author{The Stacks Project Authors}
\date{stacks-project:more-algebra.tex:local-snapshot-20261012; GNU FDL 1.2 or later}
\begin{document}
\maketitle
\setcounter{subsection}{1}
\noindent\textit{Editorial note.} All delivered mathematical text is the author's own, copied from the pinned Stacks Project chapter file more-algebra.tex; no statement or proof step was rewritten, reordered or removed.\par\medskip
\noindent\textit{Editorial note.} The chapter file pinned-original-source.tex carries only \textbackslash{}input\{preamble\} on its first line, so no preamble text lives inside the chapter. The retained file sources/stacks-more-algebra/retained-preamble.tex is a mechanical, line-for-line copy of the project's single shared preamble as staged in the snapshot, together with the project's own definitions of the reference/slogan/history display environments, which the project itself typesets as comments; nothing was added, removed or redefined by this card and no package was added.\par\medskip
\noindent\textit{Editorial note.} The source URL is the source-level Stacks Project URL, as required when more than one card is produced from one source; the per-claim tag is carried by claim\_key.\par\medskip
\noindent\textit{Editorial note.} The counted statement span includes the author's own \textbackslash{}begin\{reference\} block citing the work in which the result appears; that block is delivered verbatim and, because the project defines the reference environment as a comment environment, it produces no unresolved citation. The author's typo 'this means there is an exact sequence' region and the phrase 'we see that M\_m is a direct sum' are reproduced unchanged.\par\medskip
\noindent\textit{Allowed standard prerequisites (editorial).}\begin{itemize}
\item The characterisation of modules that are summands of direct sums of modules R/fR. invoked at proof lines 37756 and 37777. Converts the required summand statement into the surjectivity of a Hom map for all short exact sequences of the stated form.
\item Localization of modules and of Hom out of a finitely presented module. used at proof lines 37758-37772. Lets surjectivity be checked at maximal ideals and identifies the localized Hom modules.
\item The structure of finitely presented modules over a ring whose localizations satisfy the generalized valuation ring conditions. invoked at proof lines 37773-37776. Supplies the direct sum decomposition of M\_m that makes the localized map surjective.
\end{itemize}\medskip
\ifdefined\hypertarget\hypertarget{sourcepointer-1}{}\fi
\noindent\textit{Source reference.} Original lemma 00HN: focused-seeds/AGP/sources/stacks-snapshot-20261011/tags.txt, source lines 649--649. Cross-chapter navigation pointer into the Stacks Project; the standard fact it supplies is applied inside the counted proof span and the argument resting on it stays there.\par\medskip
\ifdefined\hypertarget\hypertarget{sourcepointer-2}{}\fi
\noindent\textit{Source reference.} Original lemma 0583: focused-seeds/AGP/sources/stacks-snapshot-20261011/tags.txt, source lines 6766--6766. Cross-chapter navigation pointer into the Stacks Project; the standard fact it supplies is applied inside the counted proof span and the argument resting on it stays there.\par\medskip
\ifdefined\hypertarget\hypertarget{sourcepointer-3}{}\fi
\noindent\textit{Source reference.} Original lemma 37586: sources/stacks-more-algebra/pinned-original-source.tex, source lines 37586--37586. Same-chapter navigation pointer; the statement it names is standard for the section and the argument that consumes it stays inside the counted proof span.\par\medskip
\ifdefined\hypertarget\hypertarget{sourcepointer-4}{}\fi
\noindent\textit{Source reference.} Original lemma 37638: sources/stacks-more-algebra/pinned-original-source.tex, source lines 37638--37638. Same-chapter navigation pointer; the statement it names is standard for the section and the argument that consumes it stays inside the counted proof span.\par\medskip
\ifdefined\hypertarget\hypertarget{sourcepointer-5}{}\fi
\noindent\textit{Source reference.} Original lemma 37675: sources/stacks-more-algebra/pinned-original-source.tex, source lines 37675--37675. Same-chapter navigation pointer; the statement it names is standard for the section and the argument that consumes it stays inside the counted proof span.\par\medskip
\begin{lemma}
\label{lemma-warfield}
\begin{reference}
\cite[Theorem 3]{Warfield-Decomposition}
\end{reference}
Let $R$ be a ring such that every local ring of $R$ at a maximal
ideal satisfies the equivalent conditions of
Lemma \ref{lemma-generalized-valuation-ring}.
Then every finitely presented $R$-module is a summand of a 
finite direct sum of modules of the form $R/fR$ for $f$ in $R$ varying.
\end{lemma}

\begin{proof}
Let $M$ be a finitely presented $R$-module. We first show that $M$ is a
summand of a direct sum of modules of the form $R/fR$ and at the end we
argue the direct sum can be taken to be finite. Let
$$
0 \to A \to B \to C \to 0
$$
be a short exact sequence of $R$-modules such that $fA = A \cap fB$
for all $f \in R$. By Lemma \ref{lemma-characterize-PD-modules}
we have to show that $\Hom_R(M, B) \to \Hom_R(M, C)$ is surjective.
It suffices to prove this after localization at maximal ideals
$\mathfrak m$, see
Algebra, Lemma \ref{algebra-lemma-characterize-zero-local}.
Note that the localized
sequences $0 \to A_\mathfrak m \to B_\mathfrak m \to C_\mathfrak m \to 0$
satisfy the condition that $fA_\mathfrak m = A_\mathfrak m \cap fB_\mathfrak m$
for all $f \in R_\mathfrak m$ (because we can write $f = uf'$ with
$u \in R_\mathfrak m$ a unit and $f' \in R$ and because localization
is exact). Since $M$ is finitely presented, we see that
$$
\Hom_R(M, B)_\mathfrak m = \Hom_{R_\mathfrak m}(M_\mathfrak m, B_\mathfrak m)
\quad\text{and}\quad
\Hom_R(M, C)_\mathfrak m = \Hom_{R_\mathfrak m}(M_\mathfrak m, C_\mathfrak m)
$$
by Algebra, Lemma \ref{algebra-lemma-hom-from-finitely-presented}.
The module $M_\mathfrak m$ is a finitely presented $R_\mathfrak m$-module. By
Lemma \ref{lemma-generalized-valuation-ring-modules}
we see that $M_\mathfrak m$ is a direct sum of modules
of the form $R_\mathfrak m/fR_\mathfrak m$. Thus we conclude by
Lemma \ref{lemma-characterize-PD-modules} that the map on
localizations is surjective.

\medskip\noindent
At this point we know that $M$ is a summand of
$\bigoplus_{i \in I} R/f_i R$. Consider the map
$M \to \bigoplus_{i \in I} R/f_i R$. Since $M$ is a finite $R$-module,
the image is contained in $\bigoplus_{i \in I'} R/f_i R$ for some finite
subset $I' \subset I$. This finishes the proof.
\end{proof}

\end{document}
